\section{Related Work}

\label{sec:related-work}
%\subsection{Mixed-Initiative Dialog}

\textbf{Mixed-initiative dialog}~\cite{4081977, 796083, chu-carroll-2000-mimic} refers to communication with freeflowing questions and answers from both parties. In the NLP field, the dominant chatbot paradigm adopted by large language models (LLMs) largely eschews mixed-initiative interaction: humans pose substantive questions, and the chatbot primarily responds to fulfill these requests~\cite{ouyang2022training, achiam2023gpt}. 
Recent work has sought to make dialog systems more goal-directed and proactive by incorporating mixed-initiative strategies in tasks such as creating documents~\cite{collabllm2025}, persuading users to donate to charity, enhancing users’ emotional well-being~\cite{deng2023survey, yu2023promptbasedmontecarlotreesearch, chen2023controllable, deng2024plugandplaypolicyplannerlarge}, clarifying ambiguous human requests~\cite{qian2021databasesearchresultsdisambiguation, deng2023promptingevaluatinglargelanguage, chen2024learningclarifymultiturnconversations}, or as part of an active-learning framework~\cite{Thomason2019ImprovingGN}. 
However, none of these systems addressed mixed-initiative dialog in grounded, real-world collaborative scenarios involving physical manipulation tasks.

%\subsection{Human-Robot Collaboration with Dialog}

% There are works that do HRI with dialog, but they are single initiative. There are works that do mixed-initiative HRI, but do not consider dialog.

In the human-robot interaction (HRI) field, researchers have developed \textbf{human-robot collaboration systems} that interact through language but are restricted to \textbf{single-initiative dialog}.
Some of these systems integrate LLMs as task planners or delegators~\cite{wang2024mosaicmodularassistiveinteractive, mandi2023rocodialecticmultirobotcollaboration, feng2024largelanguagemodelbasedhumanagent} for tasks like real-world cooking~\cite{wang2024mosaicmodularassistiveinteractive} and object sorting~\cite{mandi2023rocodialecticmultirobotcollaboration}.
Other systems implement a leader-follower paradigm in simulated worlds, where the leader issues natural language instructions for the follower to execute~\cite{suhr2022executinginstructionssituatedcollaborative, kojima2021continuallearninggroundedinstruction, deepmindinteractiveagentsteam2022creatingmultimodalinteractiveagents, gao2023alexaarenausercentricinteractive}.
Single-initiative HRI systems can ask humans for clarification~\cite{ren2023robotsaskhelpuncertainty} or assistance~\cite{bennetot2020artificialagentsaskhelp,6630673, 10.5555/2832747.2832901}, or inform humans of their observations~\cite{chen:jair10, 4115652, cascianelli2018full}. 
However, by supporting only single-initiative dialog, these systems lack the capacity to adapt to the evolving nature of the human, robot, and environment—limiting their capacity to find the optimal division of labor that respects user preferences~\cite{mandi2023rocodialecticmultirobotcollaboration}. 
% while also adapting to their requirements and they often have to resort to assume (unrealistically) that humans will always comply with the robot's requests~\cite{mandi2023rocodialecticmultirobotcollaboration}. 
% Other researchers have 
% While these approaches enable robots to initiate conversation in certain situations, they do not support true back-and-forth mixed-initiative dialog where humans and robots collaborate as equals to jointly decide how to best accomplish a task.
% In a related setting, researchers in Human-robot interaction (HRI) have explored a leader-follower paradigm for in simulated worlds, where the leader issues natural language instructions for the follower to execute~\cite{suhr2022executinginstructionssituatedcollaborative, kojima2021continuallearninggroundedinstruction, deepmindinteractiveagentsteam2022creatingmultimodalinteractiveagents, gao2023alexaarenausercentricinteractive}. 
% Other HRI real robot systems ask humans for help~\cite{6630673, 10.5555/2832747.2832901} or inform humans of their observations~\cite{chen:jair10, 4115652, cascianelli2018full}. However, these prior works in the HRI field are also limited to one-way communication.
% Generalizing beyond human-robot teams, natural language dialog has been studied between multi-agent robot collaborators~\cite{zhang2024efficientllmgroundingembodied, mandi2023rocodialecticmultirobotcollaboration}, where initiative and collaboration are easier to induce because the same LLM prompt can simply be provided to both agents---avoiding a lot of the challenges and unpredictability that come with adapting to the preferences and helpfulness of a human collaborator.

Some works in HRI have explored \textbf{mixed-initiative collaborative systems without dialog}, only with physical actions~\cite{4107804, natarajan2024mixedinitiativehumanrobotteamingsuboptimality, 9106686, rosero2021cooksunderstandingdynamichumanagent, paleja2024designsenablingcollaborationhumanmachine, 7379306}.
% In this paradigm, a robotic agent and a human share control of a physical task that either agent can initiate. 
% Most of the works in this paradigm have been evaluated in simulated environments~\cite{4107804, natarajan2024mixedinitiativehumanrobotteamingsuboptimality, 9106686, rosero2021cooksunderstandingdynamichumanagent, paleja2024designsenablingcollaborationhumanmachine}.
In particular, \cite{7451735} studied separate regimes of agent initiative (human-initiative, requesting help, or robot-initiative, proactively helping), but failed to support a natural human-robot dialog. 
By focusing solely on physical actions, these prior works overlook the critical role of communication in effective collaboration, thereby limiting the flexibility of the human-robot team. With \methodname{}, we enable both agents to take initiative—through both physical and verbal actions—via task-grounded dialog.
% For interested readers in mixed-initiative systems with physical actions, we refer them to \citet{7379306}, which provides surveys of them and proposes a three-dimensional taxonomy characterizing them by \textit{when}, \textit{how}, and \textit{on what aspects} of a task either agent can initiate interactions.

% Several studies explore mixed-initiative human-human~\cite{9106686, rosero2021cooksunderstandingdynamichumanagent} and human-robot interactions~\cite{paleja2024designsenablingcollaborationhumanmachine} in Overcooked-AI~\cite{carroll2020utilitylearninghumanshumanai}, a two-agent simulated collaborative cooking game, though the latter work in HRI does not consider dialog for collaboration.

% \subsection{Human-Robot Task Allocation}
Several prior works in robotics and planning have studied the problem of \textbf{human-robot optimal task allocation}, typically optimizing the time to perform a task or minimizing idle agents, posing the problem as a scheduling optimization~\cite{vats2022synergisticschedulinglearningallocation, YU2021487}.
Others have prioritized different objectives, such as safety~\cite{Faccio2024} through the formulation of a constrained optimization problem~\cite{singh2023concurrentconstrainedoptimizationunknown}.
While these solutions may achieve shorter execution times, they assume a priori known capabilities and availability of all agents, including both robots and humans.
In contrast, \methodname{} can adapt to the specific human's willingness to help by estimating its availability based on previous dialog.

% , or delegating steps of a long-horizon task between a human and a robot, where either agent can initiate dialog to express their preferences, capabilities, or to request help. 
% Several works have formulated human-robot task allocation as a scheduling problem, or as a constrained optimization problem centered on safety
% ~\citet{singh2023concurrentconstrainedoptimizationunknown} examined optimizing task allocation by building coalitions in multi-robot groups in simulation.
% ~\citet{7451735} studied 3 separate regimes of initiative: human-initiative (where it asks the robot for help), robot-initiative reactive, and robot-initiative proactive. In all regimes, the robot does not have any dialog capabilities. The authors found that a proactively helping robot yielded the best satisfaction scores. We extend this work by enabling both agents to take initiative through task-grounded dialog interactions while deciding how to delegate the steps of a long-horizon task.
